% Imporditud: tools/import_eksperimendid.py
% Allikas: Eesti füüsikaolümpiaadi eksperimendi ülesannete
%   kogu (Rael Kalda, 2023), ülesanne Ü27, lahendus L27.
\setAuthor{Erkki Tempel}
\setRound{lõppvoor}
\setYear{2018}
\setNumber{P E1}
\setDifficulty{1}
\setTopic{vedelike mehaanika}
\setVahendid{mutter, plasttops veega (0,2 l, $\rho$v = 1,0 g/cm$^3$), väike plasttops (40 ml),}
\setPunktid{10}
\setAllikas{Eksperimendi kogu Ü27, lahendus lk 44}
\setLahendusseis{toores}

\prob{Mutri tihedus}
Määrake mutri tihedus.

süstal (20 ml).

\hint

\solu
Märgime plastiktopsil veenivoo. Uputame mutri topsi. Eemaldame süstlaga topsist nii palju vett, et veetase oleks märgini - saame teada mutri ruumala Vm. Paneme väikese plastiktopsi vette ujuma ning märgime ära veenivoo. Paneme mutri väikesesse topsi nii, et ta jääb ujuma. Eemaldame süstlaga topsist nii palju vett, et veetase oleks märgini - saame teada mutri pool väljatõrjutud vee ruumala Vv ning leiame jõudude tasakaalust mutri massi mm.

mmg = $\rho$vVvg $\Rightarrow$ mm = $\rho$vVv

Teades ruumala ja massi, saame leida tiheduse $\rho$m

$\rho$m = mm = $\rho$vVv . Vm Vm
\probend
